
\subsubsection{2.1 Truthfulness Evaluation}

TruthfulQA \citep{lin2022truthfulqa} measures whether language models generate truthful answers, finding that larger models are not necessarily more truthful and may reproduce misconceptions more fluently. Subsequent work examined truthfulness in specific domains and explored interventions such as RLHF to improve truthful behavior \citep{ouyang2022instructgpt}. These studies evaluate truthfulness in isolation without examining relationships to adversarial robustness.

\subsubsection{2.2 Adversarial Robustness}

AdvGLUE \citep{wang2022advglue} provides a multi-task adversarial benchmark with word-level perturbations across GLUE tasks. Earlier work established adversarial vulnerabilities through TextFooler \citep{jin2020textfooler} and BERT-Attack \citep{li2020bertattack}. These studies characterize robustness vulnerabilities but do not examine whether robust models are also truthful.

\subsubsection{2.3 Calibration in Language Models}

Guo et al. (2017) formalized Expected Calibration Error (ECE), showing modern neural networks are often poorly calibrated. For language models, Zhao et al. (2021) demonstrated that contextual calibration improves few-shot performance. Kadavath et al. (2022) examined whether models ''know what they know,'' connecting calibration to epistemic uncertainty.

Calibration provides a plausible hypothesis for why truthfulness and robustness might correlate: a well-calibrated model accurately estimates uncertainty, potentially enabling both refusal of uncertain claims (truthfulness) and detection of anomalous inputs (robustness). This hypothesis is directly tested in this work.

\subsubsection{2.4 Multi-Dimensional Trust}

While individual trust dimensions have been studied extensively, systematic examination of their relationships is limited. Liang et al. (2023) proposed HELM for holistic evaluation across multiple dimensions, but focused on coverage rather than inter-metric correlation. This work fills this gap by quantifying the correlation between truthfulness and robustness and testing a mechanistic explanation.
